\section{Contributions}
\label{sec:conclusion:contributions}

Phase~One produces seven contributions.

\begin{enumerate}
    \item \textbf{A reconciled design specification for an integrated
    bus-network analytical platform.} The specification couples the
    AFC and GPS data streams at a documented aggregate grain, names
    six analytical themes spanning three user roles, and binds each
    theme to a chosen visualisation primitive and an interaction
    pattern. The specification is the principal output of
    Chapter~\ref{chapter:solution} and is the design contract for
    Phase~Two.

    \item \textbf{A Medallion warehouse with a materialised Gold tier
    and a dimensionally-organised Silver tier.} The warehouse design
    responds to the latency budget of the proposed interactive
    analytical surface by pre-computing seven aggregate tables in the
    Gold tier at build time. The design is defended against two
    rejected alternatives (plain Medallion with on-demand
    aggregation, and a star schema over raw views) by a measured
    benchmark against the running warehouse: the Gold tier reduces
    total chart-catalogue serving cost from approximately
    5.5~seconds per page load to approximately 0.9~seconds, a
    six-times reduction. The Silver tier's internal dimensional
    organisation is independently defended by three concrete
    row-multiplication bugs in the raw streams that the dimension
    tables resolve.

    \item \textbf{A fourteen-question stakeholder-requirements
    record.} Each analytical capability in the proposed system is
    paired with a concrete stakeholder question, a six-field
    requirements record (input aggregate, computation, output shape,
    priority tier, decision type, time-bounded success criterion),
    and a traceability mapping to a chart in the catalogue. The
    record is the source-of-truth specification against which every
    chart in the catalogue (Section~\ref{sec:solution:viz}) is
    designed.

    \item \textbf{A twenty-canvas evidence-grounded chart
    catalogue.} Each analytical theme in the catalogue is justified
    by at least one finding from the exploratory data analysis of
    Section~\ref{sec:solution:dataset}. The catalogue records the
    chosen primitive, the warehouse data source, and the dominant
    interaction pattern for every chart, and is defended against six
    rejected primitive choices. Two extensions (a city-wide
    bus-as-traffic-sensor speed grid~\cite{derevitskiy2016} and a
    per-segment bus-speed-against-posted-limit map) are designed but
    conditioned on the full fifty-day GPS extraction and on the
    ingestion of an OpenStreetMap road-speed reference, respectively;
    they are outside the twenty-canvas catalogue and recorded as
    Phase-Two extensions in Section~\ref{sec:future:extensions}.

    \item \textbf{Twelve documented data invariants for the
    integrated AFC and GPS sources.} The invariants are derived from
    two sequential passes over the data: an initial exploratory pass
    that surfaced six observations, and a subsequent systematic
    profiling pass over every raw table that surfaced an additional
    five. The invariants are recorded in
    Table~\ref{tab:dataset:quality} with their downstream effects on
    the warehouse and analytical surface. They include the calendar-
    window mismatch between the two streams, the near-disjoint
    populations of GPS and AFC drivers, the systematic
    underspecification of the route-identifier key space across
    sources, and the differential coverage of GPS instrumentation
    fields between the local sample and the full dataset.

    \item \textbf{A five-tool competitive analysis covering the
    existing artefacts that touch the Ho Chi Minh City bus network.}
    The analysis is recorded in
    Section~\ref{sec:related:tools} and establishes that no currently
    accessible tool provides a multi-stakeholder analytical surface
    over the network's historical operational and demand data. The
    one publicly accessible research demonstrator covers a single
    chart equivalent within the proposed catalogue. The five
    differentiation moats of the proposed system, school-aware demand,
    integrated AFC and GPS at the aggregate grain, in-theme stories,
    a unified multi-stakeholder surface, and a pre-tagged HCMC-specific
    anomaly catalogue, are stated explicitly.

    \item \textbf{A running prototype of the proposed system.} The
    running prototype implements the full Medallion warehouse (Bronze,
    Silver, Gold) with idempotent rebuild from the raw inputs and
    serves nine chart canvases across three of the six analytical
    themes plus one scrollytelling Story; four exploratory notebooks
    plus one programmatic profiling script produce the field-description
    tables, the exploratory figures, and the schema-map and findings
    reports referenced throughout Chapter~\ref{chapter:solution}.
\end{enumerate}
